Federal securities class actions almost always end in settlement or
dismissal, yet existing empirical work studies these outcomes in
isolation. This thesis asks how the Private Securities Litigation
Reform Act of 1995 (PSLRA), judicial geography, and observable case characteristics
jointly shape the timing and type of resolution in federal securities class actions.
It applies, to our knowledge, the first full
competing-risks survival framework to federal securities litigation,
analyzing 12,968 cases from the Federal Judicial Center's Integrated
Database (1990--2024) using cumulative incidence functions,
cause-specific Cox and Fine-Gray subdistribution models, inverse
probability of treatment weighting (IPTW), and shared frailty
models.

Three findings emerge. First, the PSLRA is consistently associated
with elevated early dismissal: a piecewise model reveals a 79\%
higher first-year dismissal hazard (HR~$= 1.789$) that decays to
null beyond two years, consistent with an early pleading-stage
filter. This association survives composition adjustment (IPTW
HR~$= 1.519$), though temporal confounding cannot be ruled out. The
settlement-side association is directionally consistent but
materially fragile.

Second, circuit geography is the dominant structural predictor of
outcomes, with settlement hazards ranging from 1.1 to 4.3 times the
Second Circuit baseline.

Third, multidistrict litigation (MDL) consolidation produces a sign
reversal visible only through competing-risks analysis---a lower
instantaneous settlement hazard coexists with higher cumulative
settlement probability---demonstrating that single-outcome
frameworks mischaracterize resolution dynamics.

Taken together, the PSLRA's strongest association is an
early-dismissal filter, while settlement suppression evidence
remains fragile.
